% =============================================================================
\section{Ejemplos de solo lectura}
% =============================================================================

\begin{frame}{Cinco nodos que solo escuchan}
  \centering\small
  \begin{tabular}{lll}
    \toprule
    Nodo & Lee & Para qué \\
    \midrule
    \cmd{telemetry\_listener} & estado, posición, GPS, actitud, batería & un resumen cada segundo \\
    \cmd{imu\_listener} & \cmd{sensor\_combined} & giroscopio y acelerómetro \\
    \cmd{gps\_status} & \cmd{vehicle\_gps\_position} & ¿ya se puede volar con GPS? \\
    \cmd{flight\_mode\_monitor} & \cmd{vehicle\_status\_v1}, \cmd{failsafe\_flags} & cada cambio de modo y failsafe \\
    \cmd{battery\_monitor} & \cmd{battery\_status\_v1} & avisar bajo un \% (y opcional RTL/LAND) \\
    \bottomrule
  \end{tabular}
  \vspace{0.6em}
  \begin{idea}
    \small No mandan nada a PX4 (salvo \cmd{battery\_monitor} si se lo piden): se pueden dejar corriendo
    \textbf{mientras} otro nodo vuela, para ver qué está pasando.
  \end{idea}
\end{frame}

\begin{frame}{El patrón: recibir en callbacks, imprimir en el timer}
  \centering
  \begin{tikzpicture}[node distance=0.35cm,every node/.style={font=\scriptsize},
      tp/.style={draw=gray,fill=gray!10,rounded corners,minimum width=3.2cm},
      mb/.style={draw=azul,fill=azulclaro,rounded corners,minimum width=1.8cm}]
    \node[tp] (t1) {vehicle\_status\_v1 (\textasciitilde5 Hz)};
    \node[tp,below=of t1] (t2) {vehicle\_local\_position (\textasciitilde50 Hz)};
    \node[tp,below=of t2] (t3) {vehicle\_attitude (\textasciitilde100 Hz)};
    \node[tp,below=of t3] (t4) {battery\_status\_v1 (1 Hz)};
    \node[mb,right=1.2cm of t1] (m1) {status\_};
    \node[mb,right=1.2cm of t2] (m2) {local\_};
    \node[mb,right=1.2cm of t3] (m3) {attitude\_};
    \node[mb,right=1.2cm of t4] (m4) {battery\_};
    \foreach \i in {1,2,3,4} \draw[-{Stealth}] (t\i) -- node[above,font=\tiny]{\cmd{= *msg}} (m\i);
    \node[draw=azul,fill=yellow!20,rounded corners,minimum height=2.6cm,right=1.2cm of m2,yshift=-0.45cm,align=center] (tm) {timer 1 s\\print\_summary()};
    \foreach \i in {1,2,3,4} \draw[-{Stealth},dashed] (m\i.east) -- (tm.west);
  \end{tikzpicture}
  \vspace{0.4em}

  \raggedright\small Cada topic llega a su ritmo; los callbacks solo guardan \textbf{la última copia}. El timer
  lee esas copias y escribe una sola línea por segundo, sin llenar la terminal.
\end{frame}

% ----------------------------------------------------------------------------- telemetry
\begin{frame}{telemetry\_listener}
  \begin{columns}[c]
    \column{0.5\textwidth}\small
    Se suscribe a 5 topics y cada segundo imprime:
    \begin{itemize}
      \item \textbf{Estado}: armado o no, y el modo (nombre y número).
      \item \textbf{Local}: x, y, altura ($-z$) y velocidad NED.
      \item \textbf{GPS}: lat, lon, alt (AMSL).
      \item \textbf{Actitud}: roll, pitch, yaw en grados (del cuaternión).
      \item \textbf{Batería}: voltaje y \%.
    \end{itemize}
    \column{0.47\textwidth}
    \begin{terminal}{ros2 run px4\_examples telemetry\_listener}
Estado: desarmado | modo: HOLD (4)\\
Local: x=0.01 y=-0.02 alt=0.03 m\\
\ \ | vel=(0.00, 0.00, 0.00) m/s\\
GPS: lat=47.3979711 lon=8.5461637\\
\ \ alt=0.1 m (AMSL)\\
Actitud: roll=0.1° pitch=-0.2° yaw=90.0°\\
Batería: 16.20 V | 100 \%
    \end{terminal}
  \end{columns}
\end{frame}

\begin{frame}[fragile]{telemetry\_listener: ¿llegaron datos?}
  \begin{lstlisting}[style=cpp,caption={telemetry\_listener.cpp}]
void print_summary() {
  // Si timestamp sigue en 0 es que todavía no ha llegado ningún mensaje
  // (los mensajes se crean con todos sus campos en 0). Suele indicar que
  // PX4 o el agente no están corriendo, o que el nombre del topic está mal.
  if (status_.timestamp == 0) {
    RCLCPP_WARN(get_logger(), "Sin datos de vehicle_status todavía (¿corre el agente?)");
    return;
  }
  ...
  \end{lstlisting}
  \small
  \begin{idea}
    \small Truco útil en cualquier nodo: un miembro con \cmd{timestamp == 0} significa ``nunca llegó un
    mensaje''. Antes de usar la posición, revisen que ya haya llegado (y que \cmd{xy\_valid} sea \cmd{true}).
  \end{idea}
\end{frame}

\begin{frame}[fragile]{telemetry\_listener: de \texttt{nav\_state} a texto}
  \begin{lstlisting}[style=cpp,caption={Una función que traduce el número a un nombre}]
const char * nav_state_name(uint8_t s) {
  switch (s) {
    case VehicleStatus::NAVIGATION_STATE_AUTO_LOITER: return "HOLD";
    case VehicleStatus::NAVIGATION_STATE_OFFBOARD:    return "OFFBOARD";
    case VehicleStatus::NAVIGATION_STATE_AUTO_LAND:   return "LAND";
    ...
    default: return "OTRO";   // cualquier otro valor
  }
}
  \end{lstlisting}
  \small \cmd{const char*} es texto fijo de C: lo que espera \cmd{\%s} en \cmd{RCLCPP\_INFO}. Cada \cmd{case}
  termina con \cmd{return}, así que no hace falta \cmd{break}.
\end{frame}

% ----------------------------------------------------------------------------- imu
\begin{frame}{imu\_listener}
  \begin{columns}[c]
    \column{0.5\textwidth}\small
    Lee \cmd{/fmu/out/sensor\_combined} (cientos de Hz) y muestra:
    \begin{itemize}
      \item \textbf{Gyro}: velocidad de giro en x, y, z (rad/s y °/s).
      \item \textbf{Accel}: aceleración en x, y, z (m/s$^2$).
      \item La frecuencia a la que llegan los mensajes.
    \end{itemize}
    Prueben girar el dron en Gazebo (o despegarlo con otro nodo) y vean cómo cambian.
    \column{0.47\textwidth}
    \begin{terminal}{ros2 run px4\_examples imu\_listener}
Gyro [rad/s]: x=+0.001 y=-0.002 z=+0.000\\
\ \ (deg/s: +0.1 -0.1 +0.0)\\
Accel [m/s²]: x=+0.012 y=-0.008\\
\ \ z=-9.806 | 250 Hz
    \end{terminal}
  \end{columns}
\end{frame}

\begin{frame}{¿Por qué el acelerómetro marca $-9.81$ quieto?}
  \begin{columns}[c]
    \column{0.45\textwidth}\centering
    \begin{tikzpicture}[scale=1]
      \fill[gray!30] (-1.6,-0.15) rectangle (1.6,0);
      \draw[thick,fill=azulclaro] (-0.6,0) rectangle (0.6,0.4);
      \draw[-{Stealth},thick,rojo] (0,0.2) -- (0,-1.1) node[below,font=\tiny] {gravedad};
      \draw[-{Stealth},thick,verde] (0.3,0.4) -- (0.3,1.5) node[above,font=\tiny] {el suelo empuja};
      \draw[-{Stealth},azul] (-1.1,0.2) -- (-1.1,-0.5) node[below,font=\tiny] {z (FRD)};
    \end{tikzpicture}
    \column{0.52\textwidth}\small
    \begin{itemize}
      \item Un acelerómetro mide la \textbf{fuerza específica}: todo menos la gravedad.
      \item Quieto en el suelo, lo único que siente es al suelo empujando \textbf{hacia arriba}.
      \item En FRD $z$ apunta hacia abajo, así que arriba es $z$ \textbf{negativa}: $(0, 0, -9.81)$.
      \item En caída libre marcaría $(0, 0, 0)$.
    \end{itemize}
  \end{columns}
\end{frame}

% ----------------------------------------------------------------------------- gps
\begin{frame}[fragile]{gps\_status}
  \begin{columns}[c]
    \column{0.5\textwidth}\small
    Lee el GPS \textbf{crudo} y dice si ya se puede volar con él. Revisa:
    \begin{itemize}
      \item \cmd{fix\_type} $\geq 3$ (3D),
      \item \cmd{satellites\_used} $\geq$ \cmd{min\_satellites} (6),
      \item \cmd{eph} $\leq$ \cmd{max\_eph} (3.0 m).
    \end{itemize}
    \begin{lstlisting}[style=bash]
ros2 run px4_examples gps_status --ros-args \
  -p min_satellites:=8 -p max_eph:=2.0
    \end{lstlisting}
    \column{0.47\textwidth}
    \begin{terminal}{ros2 run px4\_examples gps\_status}
3D | 10 sat | eph 0.9 m epv 1.4 m\\
\ \ | hdop 0.8 | 47.3979711, 8.5461637\\
\ \ GPS listo para volar
    \end{terminal}
    \footnotesize En el dron real, adentro de un edificio verán ``Todavía no:'' seguido de lo que falta.
  \end{columns}
\end{frame}

\begin{frame}{GPS crudo contra posición estimada}
  \centering
  \begin{tikzpicture}[node distance=0.5cm,every node/.style={font=\scriptsize},
      bx/.style={draw=azul,fill=azulclaro,rounded corners,minimum height=0.7cm,align=center}]
    \node[bx,fill=gray!15,draw=gray] (r) {receptor GPS};
    \node[bx,right=of r] (g) {\cmd{vehicle\_gps\_position}\\(crudo, 5--10 Hz)};
    \node[bx,right=of g] (e) {EKF2\\+ IMU, baro, mag};
    \node[bx,right=of e,fill=yellow!20] (o) {\cmd{vehicle\_global\_position}\\(estimado, rápido)};
    \draw[-{Stealth}] (r) -- (g); \draw[-{Stealth}] (g) -- (e); \draw[-{Stealth}] (e) -- (o);
    \node[below=0.2cm of g,font=\tiny] {\cmd{gps\_status}};
    \node[below=0.2cm of o,font=\tiny] {\cmd{telemetry\_listener}};
  \end{tikzpicture}
  \vspace{0.6em}

  \raggedright\small Para \textbf{volar} se usa la estimada; para \textbf{diagnosticar} el GPS (satélites,
  interferencia), la cruda.
\end{frame}

% ----------------------------------------------------------------------------- modos
\begin{frame}{flight\_mode\_monitor}
  \begin{columns}[c]
    \column{0.48\textwidth}\small
    Imprime \textbf{solo cuando algo cambia}:
    \begin{itemize}
      \item armado / desarmado,
      \item el modo (y si PX4 está en otro distinto del que pidió el usuario),
      \item failsafe activo / terminado,
      \item chequeos de pre-vuelo OK o fallando,
      \item las banderas de \cmd{failsafe\_flags} activas.
    \end{itemize}
    \column{0.5\textwidth}
    \begin{terminal}{ros2 run px4\_examples flight\_mode\_monitor}
DESARMADO\\
Modo: Hold\\
Chequeos de pre-vuelo: OK, se puede armar\\
ARMADO\\
Modo: Hold -> Offboard\\
Modo: Offboard -> Hold\\
FAILSAFE ACTIVO (modo Hold)\\
Banderas de failsafe:\\
\ \ señal offboard perdida
    \end{terminal}
  \end{columns}
\end{frame}

\begin{frame}[fragile]{flight\_mode\_monitor: comparar con lo anterior}
  \begin{lstlisting}[style=cpp,caption={flight\_mode\_monitor.cpp (recortado)}]
void on_status(const VehicleStatus & s) {
  if (first_ || s.nav_state != last_.nav_state) {
    RCLCPP_INFO(get_logger(), "Modo: %s -> %s", mode_name(last_.nav_state),
                mode_name(s.nav_state));
    // Si el modo activo no es el que pidió el usuario, PX4 lo cambió solo.
    if (s.nav_state != s.nav_state_user_intention) { ... }
  }
  last_ = s;          // guardar para la próxima comparación
  first_ = false;
}
  \end{lstlisting}
  \small Guardar el mensaje anterior y comparar es la forma de detectar \textbf{eventos} (cambios) en
  lugar de imprimir el estado una y otra vez.
\end{frame}

\begin{frame}{Úsenlo para entender qué hace PX4}
  \small Dejen \cmd{flight\_mode\_monitor} corriendo en una terminal y en otra:
  \begin{itemize}
    \item Corran \cmd{takeoff\_mission} y luego Ctrl+C: verán \textbf{Offboard $\rightarrow$ failsafe}.
    \item Corran \cmd{orbit\_command}: \textbf{Takeoff $\rightarrow$ Hold $\rightarrow$ Orbit $\rightarrow$
      Return $\rightarrow$ Land}.
    \item Cambien de modo desde QGroundControl.
  \end{itemize}
  \begin{idea}
    \small Cuando un nodo ``no hace nada'', casi siempre es porque PX4 está en un modo distinto del que
    creían. Este monitor lo dice en una línea.
  \end{idea}
\end{frame}

% ----------------------------------------------------------------------------- bateria
\begin{frame}{battery\_monitor}
  \begin{columns}[c]
    \column{0.5\textwidth}\small
    Lee la batería, avisa bajo \cmd{low\_percent} y, si se le pide, actúa:
    \begin{center}\scriptsize
      \begin{tabular}{ll}
        \toprule
        \cmd{action} & Qué hace \\
        \midrule
        \cmd{none} & solo avisa \\
        \cmd{rtl} & \cmd{NAV\_RETURN\_TO\_LAUNCH} \\
        \cmd{land} & \cmd{NAV\_LAND} ahí mismo \\
        \bottomrule
      \end{tabular}
    \end{center}
    Solo actúa si el dron \textbf{está armado}.
    \column{0.47\textwidth}
    \begin{terminal}{battery\_monitor -p low\_percent:=40.0 -p action:=rtl}
85.0 \% | 15.80 V (3.95 V/celda, 4S)\\
\ \ | -1.0 A | OK\\
BATERÍA BAJA: 39.8 \% (umbral 40 \%)\\
Regresando a casa (RTL)
    \end{terminal}
  \end{columns}
\end{frame}

\begin{frame}{El failsafe de PX4 contra el nuestro}
  \begin{columns}[c]
    \column{0.55\textwidth}\small
    PX4 ya tiene failsafe de batería (\cmd{BAT\_LOW\_THR}, \cmd{BAT\_CRIT\_THR}, \cmd{COM\_LOW\_BAT\_ACT}).
    ¿Para qué otro?
    \begin{itemize}
      \item Avisar \textbf{antes}, con un umbral propio de la misión.
      \item Que la \textbf{misión decida}: cortar, guardar dónde iba, ir a cargar (como la estación de
        carga del Nivel 3 de los retos).
    \end{itemize}
    \column{0.42\textwidth}
    \begin{ojo}
      \small En SITL \cmd{current\_a} es $-1$ (no hay sensor de corriente simulado) y la batería baja del 100 al
      50 \% en 30 s armado. No usen esos números para estimar el vuelo real.
    \end{ojo}
  \end{columns}
\end{frame}
